\documentclass{article}

% --- PAQUETES UNIFICADOS ---
\usepackage[utf8]{inputenc}
\usepackage[spanish]{babel}
\usepackage{amsmath}
\usepackage{geometry}
\usepackage{circuitikz}
\usepackage{float}
\usepackage{booktabs}
\usepackage{siunitx}
\usepackage{textcomp}
\usepackage{fancyhdr}
\usepackage[colorlinks=true, linkcolor=blue, urlcolor=blue, citecolor=blue]{hyperref}
\usepackage{hypcap} % Para que los enlaces a figuras apunten al inicio de la figura

% --- CONFIGURACIÓN GENERAL ---
\geometry{a4paper, margin=1in, top=1.2in, headheight=15pt}
\sisetup{
    output-decimal-marker = {,},
    group-separator = {.}  % ¡CORRECCIÓN: Cambiado de number-group-separator a group-separator!
}
\usetikzlibrary{babel} % Solución para conflicto circuitikz y babel
\renewcommand{\figurename}{FIG.}
\renewcommand{\listtablename}{Índice de tablas}

% --- ENCABEZADO Y PIE DE PÁGINA ---
\pagestyle{fancy}
\fancyhf{} % Limpiar encabezado y pie
\fancyhead[R]{Prácticas de electricidad}
\fancyfoot[C]{\thepage}
\renewcommand{\headrulewidth}{0.4pt}
\renewcommand{\footrulewidth}{0pt}

\begin{document}

% ==================================================================
% CONTENIDO DE LVK-01.tex
% ==================================================================
\title{PRÁCTICA 19 \\ LEY DE TENSIÓN DE KIRCHHOFF (LVK)}
\author{Prof. César Rodríguez}
\date{}
\maketitle

\tableofcontents
\newpage

\listoffigures
\newpage

\listoftables
\newpage

\section{FINALIDADES}
\begin{enumerate}
    \item Determinar por análisis o cálculo la relación entre la suma de las caídas de tensión entre los extremos de resistencias conectadas en serie y la tensión aplicada.
    \item Comprobar experimentalmente la relación determinada en 1.
\end{enumerate}

\section{INFORMACIÓN PRELIMINAR}
La solución de circuitos eléctricos complicados resulta fácil aplicando las leyes de Kirchhoff. Estas leyes fueron formuladas y publicadas por el físico Gustav Robert Kirchhoff (1824-1887), y establecieron la base del análisis moderno de redes. Son aplicables a otros circuitos que contienen una o más fuentes de tensión. En esta práctica trataremos de los que sólo tienen una fuente de tensión.

\subsection{Ley de tensión}
Consideremos el circuito serie sencillo de la figura 19-1. Hemos establecido antes que las resistencias conectadas en serie $R_1, R_2, R_3$ y $R_4$ se pueden reemplazar por su resistencia total o equivalente $R_T$ sin que resulte afectada la corriente $I_T$ en el circuito. La corriente $I_T$ puede pues ser calculada por la ley de Ohm.

\begin{equation}
    I_T = \frac{V}{R_T}
\end{equation}

También se deduce de la ecuación (19-1) que
\begin{equation}
    V = I_T \times R_T
\end{equation}

En el circuito conectado en serie de la figura 19-1
\begin{equation}
    R_T = R_1 + R_2 + R_3 + R_4
\end{equation}

\begin{figure}[H]
    \centering
    \begin{circuitikz}[american]
        \draw (0,0)
        to[battery, l=$V$, invert] (0,3)
        to[short, i>^={$I_T$}] (2,3)
        to[R, l_=$R_1$, v^=$V_1$] (4,3)
        to[R, l_=$R_2$, v^=$V_2$] (6,3)
        to[R, l_=$R_3$, v^=$V_3$] (8,3)
        to[R, l_=$R_4$, v^=$V_4$] (10,3)
        to[short] (10,0)
        to[short] (0,0);
    \end{circuitikz}
    \caption{--- $V = V_1 + V_2 + V_3 + V_4$.}
    \label{fig:19-1}
\end{figure}

Sustituyendo el valor equivalente $R_T$ de la ecuación (19-3) en la ecuación (19-2) se tiene
\[
    V = I_T(R_1 + R_2 + R_3 + R_4)
\]
o
\begin{equation}
    V = I_T \times R_1 + I_T \times R_2 + I_T \times R_3 + I_T \times R_4
\end{equation}

Una de las características de un circuito serie es que la corriente es la misma en cualquier punto del circuito. $I_T$ es pues la corriente total en $R_1, R_2, R_3$ y $R_4$. De esto se deduce que
\begin{align*}
    I_T \times R_1 &= V_1, \quad \text{caída de tensión entre los extremos de } R_1 \\
    I_T \times R_2 &= V_2, \quad \text{caída de tensión entre los extremos de } R_2 \\
    I_T \times R_3 &= V_3, \quad \text{caída de tensión entre los extremos de } R_3 \\
    I_T \times R_4 &= V_4, \quad \text{caída de tensión entre los extremos de } R_4
\end{align*}

Ahora se puede escribir la ecuación (19-4) como sigue:
\begin{equation}
    V = V_1 + V_2 + V_3 + V_4
\end{equation}

La ecuación (19-5) es la expresión matemática de la ley de tensión de Kirchhoff para resistencias conectadas en serie en un circuito cerrado.

Es evidente que se puede generalizar la ec. (19-5) para circuitos que contienen una o más resistencias conectadas en serie en

% \newpage

% ==================================================================
% CONTENIDO DE LVK-02.tex
% ==================================================================

\begin{figure}[H]
    \centering
    \begin{circuitikz}[american]
        % Main path: Battery -> R1 -> Parallel Block 1 (R2, R3) -> R4 -> Parallel Block 2 (R6, R7) -> R5 -> Battery
        \draw (0,0) to[battery, l=$V$, invert] (0,3) % Batería
            to[R, l_=$R_1$, v^=$V_1$] (2,3) node[right] {A} % R1 y nodo A
            coordinate (A_end); % Coordenada para el final de R1 / inicio del bloque paralelo 1

        % Bloque paralelo 1 (R2, R3) - Horizontal
        \coordinate (B_start) at (4,3); % Coordenada para el final del bloque paralelo 1 / inicio de R4
        \draw (A_end) -- (A_end |- 0,4) to[R, l_=$R_2$] (B_start |- 0,4) -- (B_start); % Rama superior (R2)
        \draw (A_end) -- (A_end |- 0,2) to[R, l_=$R_3$] (B_start |- 0,2) -- (B_start); % Rama inferior (R3)
        \node[left] at (B_start) {B}; % Nodo B

        % R4
        \draw (B_start)
            to[R, l_=$R_4$, v^=$V_3$] (6,3) node[right] {C} % R4 y nodo C
            coordinate (C_end); % Coordenada para el final de R4 / inicio del bloque paralelo 2

        % Bloque paralelo 2 (R6, R7) - Horizontal
        \coordinate (D_start) at (8,3); % Coordenada para el final del bloque paralelo 2 / inicio de R5
        \draw (C_end) -- (C_end |- 0,4) to[R, l_=$R_6$] (D_start |- 0,4) -- (D_start); % Rama superior (R6)
        \draw (C_end) -- (C_end |- 0,2) to[R, l_=$R_7$] (D_start |- 0,2) -- (D_start); % Rama inferior (R7)
        \node[left] at (D_start) {D}; % Nodo D

        % R5 y retorno a la batería
        \draw (D_start)
            to[R, l_=$R_5$, v^=$V_5$] (10,3) % R5
            -- (10,0) to[short, i_>=$I_T$] (0,0); % Cierre del circuito

        % Etiquetas de voltaje V2, V4
        \draw[<->] (A_end |- 0,5) -- (B_start |- 0,5) node[midway, above] {$V_2$};
        \draw[<->] (C_end |- 0,5) -- (D_start |- 0,5) node[midway, above] {$V_4$};
    \end{circuitikz}
    \caption{Aplicación de ley de tensión de Kirchhoff a un circuito serie-paralelo.}
    \label{fig:19-2}
\end{figure}

¿Es consistente esta ecuación con la ec. (19-5)? Sí, porque pasando los términos del segundo miembro de la ecuación (19-5) al primer miembro, tendremos \( V - V_1 - V_2 - V_3 - V_4 = 0 \), cuyo resultado es idéntico al de la ecuación (19-6).

\begin{figure}[H]
    \centering
    \begin{circuitikz}[american]
        \draw (0,0) node[below] {E}
            to[battery, l^=$V$, v_<=$V_5$, invert] (0,3) node[above] {F}
            to[R, l_=$R_1$, v^>=$V_1$] (3,3) node[above] {A}
            to[R, l_=$R_2$, v^>=$V_2$] (6,3) node[above] {B}
            to[R, l_=$R_3$, v^>=$V_3$] (9,3) node[above] {C}
            to[R, l_=$R_4$, v^>=$V_4$] (12,3) node[above] {D}
            to[short] (12,0)
            to[short, i_>={corriente}] (0,0);
    \end{circuitikz}
    \caption{Asignación convencional de polaridad a las tensiones en un circuito cerrado.}
    \label{fig:19-3}
\end{figure}

% \newpage

% ==================================================================
% CONTENIDO DE LVK-03.tex
% ==================================================================

\section{RESUMEN}
La ley de Kirchhoff se puede expresar de dos maneras:
\begin{enumerate}
    \item La suma de caídas de tensión en un circuito cerrado es igual a la tensión aplicada; o
    \item La suma algebraica de las tensiones en un circuito cerrado es igual a cero.
\end{enumerate}

\section{AUTOEXAMEN}
\begin{enumerate}
    \item En la figura 19-1, $V_1 = \SI{10}{V}, V_2 = \SI{12}{V}, V_3 = \SI{20}{V}, V_4 = \SI{15}{V}$. La tensión aplicada $V$ debe ser pues igual a \dotfill{} V.
    \item En la figura 19-2, $V_1 = \SI{15}{V}, V_2 = \SI{20}{V}, V_4 = \SI{27}{V}, V_5 = \SI{9}{V}$ y $V = \SI{100}{V}$. La tensión $V_3 = $ \dotfill{} V.
\end{enumerate}

\section{MATERIALES NECESARIOS}
Fuente de alimentación: Fuente de corriente continua variable regulada. \\
Equipo: EVM, miliamperímetro 0-10. \\
Resistencias: \textonehalf{} W \SI{330}{\ohm}, \SI{470}{\ohm}, \SI{1200}{\ohm}, \SI{2200}{\ohm}, \SI{3300}{\ohm}, \SI{4700}{\ohm}, \SI{5600}{\ohm} y \SI{10000}{\ohm}. \\
Diversos: Interruptor unipolar.

\section{PROCEDIMIENTO}
\subsection{Ley de tensión}
\begin{enumerate}
    \item Conectar el circuito de la figura 19-1 utilizando los valores $R_1, R_2, R_3$ y $R_4$ consignados en la tabla 19-1.
    \item Ajustar la salida de la fuente de alimentación para que $V = \SI{40}{V}$. Medir y anotar esta tensión en la tabla 19-2. Medir y anotar las tensiones $V_1, V_2, V_3$ y $V_4$ entre los extremos de $R_1, R_2, R_3$ y $R_4$, respectivamente. Sumar $V_1, V_2, V_3$ y $V_4$ y anotar la suma en la columna correspondiente de la tabla 19-2.
    \item Conectar el circuito de la figura 19-2 utilizando los valores $R_1, R_2$, etc., consignados en la tabla 19-1.
    \item Ajustar la salida de la fuente de alimentación para $V = \SI{40}{V}$. Medir esta tensión y anotarla en la tabla 19-2. Medir y anotar las tensiones $V_1, V_2, V_3, V_4, V_5$ como indica la figura 19-2. Sumar $V_1, V_2, V_3$, etc., y anotar el resultado en la columna correspondiente.
\end{enumerate}

\subsection{Problema de diseño (adicional)}
\begin{enumerate}
    \setcounter{enumi}{4}
    \item Diseñar un circuito serie-paralelo como el de la figura 19-2 utilizando las resistencias mencionadas en esta práctica, con una tensión aplicada de \SI{35}{V}, que absorban una corriente total de \SI{5}{mA}, aproximadamente. Dibujar el esquema del circuito, incluyendo los valores medidos de todas las resistencias. Presentar todos los cálculos.
    \item Conectar el circuito y ajustar la tensión de la fuente de alimentación a \SI{35}{V}. Medir la corriente $I_T$ e indicar este valor en el esquema.
\end{enumerate}

\section{PREGUNTAS}
\begin{enumerate}
    \item ¿Cómo es la suma de $V_1, V_2, V_3$ y $V_4$ (apartado 2) de la tabla 19-2 comparada con $V$? Explicar la diferencia, si la hay.
    \item ¿Cómo es la suma de $V_1, V_2 \dots V_5$ (apartado 4) de la tabla 19-2 comparada con $V$? Explicar la diferencia, si la hay.
\end{enumerate}

\newpage

\begin{table}[H]
    \centering
    \caption{Valores nominales de las resistencias.}
    \label{tab:19-1}
    \begin{tabular*}{\textwidth}{@{\extracolsep{\fill}}lcccccccc}
    \toprule
     & $R_1$ & $R_2$ & $R_3$ & $R_4$ & $R_5$ & $R_6$ & $R_7$ & $R_8$ \\
    \midrule
    Valor nominal, (\si{\ohm}) & \num{330} & \num{470} & \num{1200} & \num{2200} & \num{3300} & \num{4700} & \num{10000} & \num{5600} \\
    \bottomrule
    \end{tabular*}
\end{table}

\vspace{1cm}

\begin{table}[H]
    \centering
    \caption{Mediciones de tensión en los circuitos.}
    \label{tab:19-2}
    \begin{tabular*}{\textwidth}{@{\extracolsep{\fill}}lccccccc}
    \toprule
    Operación & $V$ & $V_1$ & $V_2$ & $V_3$ & $V_4$ & $V_5$ & $V_1+V_2+V_3+V_4+V_5$ \\
    \midrule
    2 & & & & & & X & \\
    4 & & & & & & & \\
    \bottomrule
    \end{tabular*}
\end{table}

% \newpage

% ==================================================================
% CONTENIDO DE LVK-04.tex
% ==================================================================

\subsection{Preguntas adicionales}
\begin{enumerate}
    \setcounter{enumi}{2} % Empezar la numeración en 3
    \item Enunciar la ley de tensión de Kirchhoff. Escribir también la ecuación correspondiente.
    
    \item Explicar los resultados del problema de diseño, apartado 6.
    
    \item ¿Por qué no sería posible diseñar un circuito que diera exactamente una corriente de \SI{5}{mA}?
\end{enumerate}

\section{Respuestas a las preguntas de auto-examen}
\begin{enumerate}
    \item 57.
    \item 29.
\end{enumerate}

\end{document}